%! Unit 2: Functions, Transformations, and Graph Analysis — Unit Packet Cover (KEY, 2pp)
%
% Page 1 is byte-identical to the student cover: both wrappers \input the same
% unit_cover/body.tex, so the two packets cannot drift.
%
% Page 2 is TEACHER-ONLY and is the reason this file exists. Per the teachernote
% convention, a unit test's answer rationale and Part D scoring must NOT sit at
% the foot of practice_test_key/actual_test_key: the practice test is bound into
% the STUDENT packet, so its rationale would ride along in front of students, and
% a note in a key also makes the key longer than its blank. shared/unit.mk merges
% this file into the KEY packet only. Keep the notes to ONE page — cover + notes
% is a single double-sided sheet.
\documentclass[10pt]{article}
\usepackage{atda-article}
\usepackage{atda-boxes}
\usepackage{ltablex}
\keepXColumns

\begin{document}

%! Unit 2: Functions, Transformations, and Graph Analysis — the unit packet cover sheet.
%
% THIS FILE IS THE COVER. Both wrappers \input it — unit_cover/main.tex (student
% packet) and unit_cover_key/main.tex (key packet, which adds a page 2 of exam
% scoring notes) — so page 1 can never drift between the two packets. Edit the
% cover here, never in a wrapper.

% Full-bleed royal banner
\begin{tikzpicture}[remember picture, overlay]
  \fill[royal] (current page.north west)
    rectangle ([yshift=-0.9in]current page.north east);
\end{tikzpicture}
\vspace{-0.8in}

\begin{center}
  {\color{white}\textbf{\LARGE Algebra, Trigonometry, and Data Analysis}}\\[4pt]
  {\color{mistmid}\large\textbf{Unit 2: Functions, Transformations, and Graph Analysis}}\\[6pt]
  {\color{white}\textbf{\large Unit Overview \quad\textbar\quad Eight Lessons \quad\textbar\quad Unit Test}}
\end{center}

\vspace{0.20in}
\namedateperiod
\vspace{0.14in}

\noindent
\small
In Unit 1 you were handed an expression and asked to rewrite it. In Unit 2 you are handed a
\textbf{graph} and asked what it says. That is not a change of style --- it is what this course
asks for: in AFDA the characteristics of a function are investigated \emph{from the graph}. So
every lesson here is one more question you can ask a picture. What are its inputs and outputs?
What moved it off its parent? Where does it cross, turn, level off, or change rules? And, at the
end: given a table of numbers and no picture at all, \emph{which family does this even belong to?}

\vspace{0.12in}

% BOXGUARD: all three boxes below are breakable and all three fit on page 1 today
% with room to spare, so these guards are inert as written. They are sized to
% their own boxes so that if the cover ever grows, a box moves whole to a second
% page instead of splitting and stranding a stub — the guard is the cheap
% insurance, and each was verified not to cost this sheet its single page.
\boxguard[14]
\begin{tocbox}
\small
\renewcommand{\arraystretch}{1.32}
\begin{tabularx}{\linewidth}{@{} c >{\raggedright\arraybackslash}p{5.0cm} l X @{}}
\rowcolor{mist}
\textbf{\#} & \textbf{Lesson} & \textbf{SOL} & \textbf{The one idea} \\
\midrule
2.0 & Unit Opener                       & ---                       & A graph is read, not solved \\
2.1 & Notation, Domain, and Range       & \texttt{AF.2a, e}         & $f(3)=10$ and the point $(3,10)$ are one sentence \\
2.2 & Parent Functions \& Transformations & \texttt{AF.1a--b}       & Outside the parentheses, outputs; inside, inputs \\
2.3 & Equations and Graphs Both Ways    & \texttt{AF.1d, f}         & A matching picture is not proof \\
2.4 & Intercepts, Zeros, and Extrema    & \texttt{AF.2b--d}         & Say which axis the number came from \\
2.5 & End Behavior and Asymptotes       & \texttt{AF.2f--g}         & Getting closer and closer is not arriving \\
2.6 & Piecewise-Defined Functions       & \texttt{AF.2}             & Which piece owns this input? \\
2.7 & Choosing and Comparing Models     & \texttt{AF.1c, e, g; 2h}  & Chosen by how outputs \emph{change}, not by shape \\
\end{tabularx}
\end{tocbox}

\vspace{0.12in}

\boxguard[16]
\begin{spiralbox}
\small
\begin{enumerate}[leftmargin=1.5em, itemsep=3pt, topsep=2pt]
  \item \textbf{Before you write a number down, say which axis it came from.} An interval of
        \emph{inputs} and a stretch of \emph{outputs} can be the same two numbers, in the same
        units, and mean nothing alike. This is also why an extreme is reported as a
        \emph{value and a location} --- ``\$1800'' is half an answer.
  \item \textbf{Every transformation answers one question: did the outputs change, or did the
        inputs?} Outside the parentheses touches outputs; inside touches inputs. The horizontal
        rule is never ``the sign flips'' --- it is substitution: $f(x-3)$ at $x=5$ computes $f(2)$.
  \item \textbf{Getting closer and closer is not arriving.} That is why ``none'' is so often the
        right answer to ``what is the zero / the maximum?'' And an asymptote is a \emph{line}, so
        it gets an \emph{equation}: $y=20$, not $20$.
  \item \textbf{A rule owns only its own stretch of the domain.} A piecewise function is one
        function, not several; ask ``which piece owns this input?'' before anything else. A
        \emph{corner} is a change of rule, not a hole --- it is not a break.
  \item \textbf{A family is chosen by how the outputs \emph{change}, not by what the graph looks
        like.} ``It curves upward'' rules out linear and nothing else. Constant first differences
        mean linear, constant \emph{second} differences mean quadratic, a constant ratio means
        exponential --- over equal steps in $x$.
\end{enumerate}
\end{spiralbox}

\vspace{0.12in}

\boxguard[8]
\begin{remindbox}
\small
\textbf{The unit test has four parts and is worth 100 points:} Part A vocabulary matching
(8 pts), Part B multiple choice (12 pts), Part C short answer and computation (50 pts), and
Part D extended response (30 pts). A \textbf{practice test} with the same structure and
different numbers is bound at the back of this packet --- work it with no notes, then check it
against the key. Most of this test is a \emph{graph you have never seen}, so practice reading
one out loud: domain, range, intercepts, zeros, rising and falling, extremes with their
locations, and what happens at the far ends.
\end{remindbox}


\newpage

\begin{headlinebox}{royal}{\color{white}\bfseries Unit 2 --- Exam Scoring Notes (Teacher Copy)}\end{headlinebox}

\vspace{4pt}
\noindent\footnotesize
Both unit assessments are parallel forms: same blueprint, same ideas, different numbers, with the
vocabulary definitions reworded and relettered \emph{and} the multiple-choice options reordered, so
no answer letter transfers between forms. \textbf{P} = practice test (bound in the student packet),
\textbf{A} = actual test (distributed at test time only). 100 points each.

\vspace{4pt}

\boxguard[10]
\begin{teachernote}[Part A --- Vocabulary (8 pts, 1 each)]
\footnotesize
\textbf{P:} 1--E, 2--A, 3--C, 4--H, 5--I, 6--J, 7--B, 8--F. \qquad
\textbf{A:} 1--C, 2--D, 3--F, 4--G, 5--J, 6--H, 7--I, 8--B.\\[2pt]
Both forms carry \textbf{two distractor definitions} that match no term (P: D and G; A: A and E) ---
a \emph{dilation} and a \emph{$y$-intercept}. They are there so a student who knows six terms cannot
back into the last two by elimination. Do not give part credit for confusing \emph{translation} with
\emph{dilation}: ``did the shape change?'' is the whole distinction, and it is one sentence at review.
The pairing to watch is \textbf{domain/range} --- students who miss both have swapped them, which is
one error, not two.
\end{teachernote}

\boxguard[10]
\begin{teachernote}[Part B --- Multiple Choice (12 pts, 2 each)]
\footnotesize
\textbf{P:} 1--B, 2--C, 3--D, 4--B, 5--C, 6--D. \qquad \textbf{A:} 1--A, 2--B, 3--A, 4--C, 5--A, 6--B.\\[2pt]
\textbf{Each of the six items is one of the unit's six error classes, in lesson order}, so the miss
pattern is a diagnosis, not a score. Item 1 is flipped notation (2.1) --- $f(3)=10$ read as $f(10)=3$.
Item 2 is the backwards horizontal shift (2.2); refuse ``the sign flips'' as an explanation at review
and replace it with a substitution. \textbf{Item 3 is the axis error and it is the one to reteach} ---
every option is in dollars, so units cannot settle it; the only fix is naming the axis. Item 4 is
\emph{approaches is not reaches} (2.5). Item 5 is \emph{a corner is not a break} (2.6). Item 6 is
2.7's whole spine: ``it curves upward'' rules out linear and nothing else, and option (A)/(C) is the
student who stopped at the picture.
\end{teachernote}

\boxguard[10]
\begin{teachernote}[Part C --- Short Answer \& Computation (50 pts, 5 each)]
\footnotesize
One item per content lesson, in lesson order, then three synthesis items: C1 notation, domain and
range (2.1), C2 transformations and range (2.2), C3 equation from a vertex and a point, verified
(2.3), C4 the full graph read (2.4), C5 asymptote and end behavior (2.5), C6 piecewise (2.6),
C7 choose the family from a table (2.7), C8 interval notation and a restricted domain, C9 the same
question answered algebraically \emph{and} graphically (\texttt{AF.1e}), C10 the characteristics
checklist off an equation. \textbf{Suggested split: 3 pts method, 2 pts answer.}\\[2pt]
\textbf{Three deductions to apply consistently.} (i) On \textbf{C4(d)} an extreme reported as a value
with \emph{no location} is 3 of 5 --- \texttt{AF.2c} names ``the location and value,'' and this is the
unit's most repeated instruction. (ii) On \textbf{C5} a horizontal asymptote written as a bare number
($90$ / $500$) rather than an \emph{equation} ($y=90$ / $y=500$) loses 1; and (c) and (d) must carry
\emph{different} reasons --- one reason given twice is 3 of 5. (iii) On \textbf{C7} a correct family
name with no differences or ratio shown earns 2 of 5: the justification \emph{is} the item.
On \textbf{C2(c)} note the range of $-x^2$ closes at $0$ (a bracket) while the range of $-2^x$ does
not (a parenthesis) --- the asymptote is the reason, and it is worth one sentence at review.
\end{teachernote}

\boxguard[10]
\begin{teachernote}[Part D --- Extended Response (30 pts, 15 each)]
\footnotesize
\textbf{D1 (\texttt{AF.2a--d, f}) --- 15 pts:} 2 pts domain and range; 3 pts both zeros \emph{with}
the break-even meaning; 3 pts the three intervals; \textbf{4 pts part (d)}; 3 pts part (e).
\textbf{Part (d) is the assessed idea and it is deliberately not a point}: the maximum sits on a flat
stretch, so the location is the \emph{interval} $[7,10]$ (P) / $[6,9]$ (A). A student naming one
endpoint has the value right and the location half right --- 2 of those 4. In part (e), full credit
needs a price \emph{and} a reason read off the graph; ``because it's the highest point'' with no
price named earns 1.\\[2pt]
\textbf{D2 (\texttt{AF.1c, g}; \texttt{AF.2h}; \texttt{DA.1d}) --- 15 pts:} 5 pts all three rows
computed with the constant one named; 4 pts the model and the prediction ($1944$ shares, P /
$5120$ students, A); \textbf{3 pts part (c)}; 3 pts part (d). \textbf{Part (c) is the item the unit
was built toward}: the student's arithmetic is correct and the conclusion is not, so full credit
requires saying that growing jumps rule out \emph{linear} and nothing more, and that quadratic needs
constant \emph{second} differences. A response that only re-states the right answer (``it's
exponential'') without addressing what the other student got wrong earns 1. On part (d), the word
\emph{extrapolation} must be used and the point must be that the model is being read outside the
window it was built from --- ``the model is wrong'' is the misconception, not the answer.
\end{teachernote}

\end{document}
